\section{The continuation problem and the precise scope}
\label{sec:context}

\subsection{What is already available}
The canonical manuscript connects nested sums, Gamma identities, Hurwitz
zeta jets, and generalized Stieltjes constants. Its preface carefully
separates ordinary mathematical proofs, exact finite algebra, and numerical
evidence \cite{repo:main}. That distinction is particularly relevant here:
a formal rank statement about coefficient matrices does not assert
arithmetic independence of the associated special values.

The preserved \emph{Determinantal Harmonic Resonance} report evaluates a
Gauss--Vandermonde partition function and treats its first logarithmic
resonance. For critical rational weights with simple poles on a common
integer grid, its finite moment matrix contains harmonic dependence in
only one entry. This gives depth-two cutoff generators even when the
number of summation variables is arbitrary \cite{repo:dhr}. Its further
research section asks for a classification with repeated poles, a rank-based
depth theorem, and extensions to several residue classes. Those are the
specific source questions addressed in this article.

Some nearby questions have already moved beyond that report. The current
incoming archive \emph{Polygamma Reconstruction and Resonance} records a
solution of the next two-particle resonance and of higher spatial Mellin
coefficient questions. The current \emph{Harmonic Tails, Shifts and Twists}
archive records the gap-four tail-rank family and further shifted identities
\cite{repo:incoming}. We do not claim those results again. Our principal
spectral deformation is a finite difference of an independently variable
power; it is not the same deformation as a moving denominator parameter
in the earlier Gauss ensemble.

\subsection{The three main advances}
The algebraic classification in Theorems~\ref{thm:classification}
and~\ref{thm:critical} is a necessary and sufficient test, not merely a
construction of examples. At nonzero critical decay, it identifies exactly
when repeated-pole harmonic data remain rank one. Its finite-difference
interpretation explains why arbitrarily high pole multiplicities need not
increase the depth of a cutoff generator.

The spectral identity in Theorem~\ref{thm:single} evaluates a new explicit
family for every particle number and every complex order in its convergence
domain. The polarized version, Theorem~\ref{thm:polarized}, permits independent
orders and independent shifts. Despite this freedom, only single Hurwitz
zeta functions appear: no product of two such functions is needed.

The normal subtraction in Theorem~\ref{thm:subtraction} supplies ordinary
convergent series for every associated Stieltjes jet. Its pole-cancelling
factor is explicit. The hard-cutoff constant and spectral finite part are
then compared in Corollary~\ref{cor:cutoff}; their difference is not discarded.

\subsection{Classical ingredients and novelty qualification}
Finite differences, Cauchy--Binet, Gamma differentiation, and the Hurwitz
shift formula are classical. Rational summability and discrete residues
have an extensive theory, including the work of Chen--Singer and
Arreche--Sitaula \cite{chen-singer,arreche-sitaula}. We do not present the
basic discrete-residue criterion as new. Our use of repeated differences
is tied to a prescribed finite Hankel matrix and its critical decay:
this gives the explicit harmonic-pencil classification and the special-value
identities developed below.

The new statements are research contributions relative to the inspected
repository material. The source inspection and literature checks do not
establish global historical priority. Nor do they establish that the
resulting polylogarithmic representations have minimal arithmetic depth.
The text makes no claim of proof-assistant formalization.

\subsection{Conventions}
We fix an integer $r\ge2$ and write
\begin{equation}\label{eq:d}
 d=2r-2,\qquad \Delta f(x)=f(x+1)-f(x).
\end{equation}
Thus $d$ is even. A sum over $n_1<\cdots<n_r$ is ordered; it has no
additional factor $r!$. Put
\[
 \Delta(\bm n)=\prod_{1\le i<j\le r}(n_j-n_i).
\]
The repeated use of $\Delta$ for a finite-difference operator and a
Vandermonde is distinguished by its argument.

Throughout the spectral sections, $\Rea a>0$ and
$(n+a)^{-s}=\exp(-s\log(n+a))$ uses the principal logarithm. All
parameter derivatives are derivatives of functions on open complex
domains, never derivatives inferred from a formula known only at integer
orders. We write
\[
 (s)_k=s(s+1)\cdots(s+k-1),\qquad
 H_N^{(q)}=\sum_{j=1}^N j^{-q},\qquad H_0^{(q)}=0.
\]
The generalized Stieltjes convention is
\begin{equation}\label{eq:stieltjes}
 \zeta(1+u,a)=\frac1u+
 \sum_{m\ge0}\frac{(-1)^m}{m!}\gamma_m(a)u^m.
\end{equation}
In particular $\gamma_0(a)=-\psi(a)$ and $\gamma_0(1)=\gamma$.
A superscript on $\zeta$ differentiates its first variable unless an
argument derivative is explicitly indicated. We use the standard Hurwitz
and Gamma identities with these conventions \cite{dlmf:hurwitz,dlmf:gamma}.

A matrix pencil is a matrix whose entries are linear forms in independent
formal variables. Its rank means rank over the rational-function field in
those variables. This is not the rank of a matrix obtained by treating
numerical harmonic values as independent numbers.
